\apendice{Plan de Proyecto Software}

\section{Introducción}
Esta sección se detalla la gestión y planificación del proyecto
\nombreProyecto. Se aborda la metodología de trabajo seguida, la
temporalización de las tareas a lo largo de los meses de desarrollo, la gestión
de riesgos y, finalmente, un estudio de viabilidad que engloba tanto los
aspectos económicos como las consideraciones legales del software desarrollado.

\section{Metodología y planificación temporal}
Para la organización del trabajo, se ha optado por seguir un marco de trabajo
ágil adaptado a las necesidades de un proyecto individual académico, basándose
en los principios fundamentales de \textit{Scrum}. El desarrollo ha abarcado un
periodo total de casi ocho meses, dando comienzo el 16 de octubre de 2025 y
finalizando con su presentación en junio de 2026.

El ciclo de vida del desarrollo se estructuró a alto nivel en seis grandes
fases consecutivas:
\begin{enumerate}
    \item \textbf{Investigación y Viabilidad:} Análisis del estado del arte y selección tecnológica.
    \item \textbf{Módulo de Ingesta:} Desarrollo de los adaptadores para extraer datos de \href{https://github.com}{GitHub}.
    \item \textbf{Procesamiento \acrshort{nlp} y Modelado:} Fase central de experimentación con modelos de temas.
    \item \textbf{Interfaz Web:} Desarrollo de los cuadros de mando interactivos en Streamlit.
    \item \textbf{Pruebas y Despliegue:} Refactorización, creación de tests y contenerización Docker.
    \item \textbf{Redacción de la Memoria:} Documentación técnica y académica en \LaTeX.
\end{enumerate}

A nivel operativo, el proyecto se dividió en iteraciones o \textit{Sprints} de
entre dos y tres semanas de duración, ajustando este margen de flexibilidad en
función de la complejidad del módulo a desarrollar en cada momento. Las
reuniones periódicas de seguimiento con los tutores del proyecto ejercieron la
doble función de \textit{Sprint Planning} (para priorizar las siguientes
tareas) y \textit{Sprint Review} (para evaluar los incrementos de software
logrados).

Para el seguimiento de las tareas y el control de versiones, se utilizó la
plataforma de gestión ágil \textbf{Zube}, la cual se integra nativamente de
forma bidireccional con las \textit{issues} del repositorio en
\href{https://github.com/zcc1001/tfg-topics}{GitHub}. Esta herramienta permitió
gestionar el \textit{Product Backlog}, asignar prioridades y trazar el progreso
visual de cada tarea mediante un tablero Kanban a lo largo de los diferentes
\textit{Sprints}, tal y como se ilustra en la Figura
\ref{fig:plan/zube_sprints.png}.

\imagen{plan/zube_sprints.png}{Vista del panel de gestión de Sprints en la plataforma Zube}

Esta integración bidireccional garantiza que cualquier actualización en el
estado de una tarea dentro de la interfaz de Zube se sincronice de forma
inmediata con el repositorio del proyecto. En particular, la gestión del
desarrollo se apoyó en los siguientes elementos clave:
\begin{itemize}
    \item \textbf{Gestión de incidencias y tareas (\textit{Issues}):} Se emplearon las \href{https://github.com/zcc1001/tfg-topics/issues}{\textit{issues} de GitHub} (Figura~\ref{fig:plan/github_issues.png}) para documentar individualmente cada tarea de desarrollo, corrección de errores o mejora, sirviendo como el elemento atómico de trabajo y discusión.
    \item \textbf{Planificación de hitos (\textit{Milestones}):} El avance temporal del proyecto se organizó en hitos dentro de \href{https://github.com/zcc1001/tfg-topics/milestones}{GitHub} (Figura~\ref{fig:plan/github_milestrones.png}), agrupando las tareas planificadas para cada \textit{sprint} y marcando fechas de entrega específicas.
    \item \textbf{Tablero Kanban:} A través de la pestaña principal de Zube (Figura \ref{fig:plan/zube_tab.png}), se dispuso de un tablero Kanban interactivo para visualizar de forma clara el flujo de trabajo (desde la bandeja de entrada hasta la validación y finalización de tareas), optimizando la distribución del esfuerzo.
    \item \textbf{Gráfico de trabajo restante (\textit{Burndown Chart}):} Para monitorizar el ritmo del desarrollo durante cada iteración y asegurar el cumplimiento de los objetivos fijados, se utilizaron los gráficos de \textit{burndown} generados por Zube (Figura~\ref{fig:plan/zube_burndown.png}), los cuales muestran la tendencia diaria de cierre de tareas.
\end{itemize}

\imagen{plan/github_issues.png}{Panel de gestión de Issues en GitHub}
\imagen{plan/github_milestrones.png}{Vista de los Milestones planificados en GitHub}
\imagen{plan/zube_tab.png}{Pestaña de gestión de tareas y tablero Kanban en Zube}
\imagen{plan/zube_burndown.png}{Gráfico Burndown de seguimiento del trabajo en Zube}

\subsection{Evolución real del proyecto por sprints}

Además de la planificación inicial, el seguimiento iterativo del desarrollo se
gestionó mediante \textit{sprints} cerrados en la plataforma Zube y en
\href{https://github.com/zcc1001/tfg-topics}{GitHub}. Esta trazabilidad permite
contrastar el plan previsto con la evolución real del trabajo y evidenciar el
grado de cumplimiento.

En total se completaron doce iteraciones (Sprint 0 a Sprint 13), con cierre
íntegro de tareas en todos los casos. La evolución del proyecto siguió una
secuencia progresiva: arranque y viabilidad inicial, construcción del conjunto
de datos, consolidación de los módulos de ingesta y procesamiento, mejoras
funcionales y visuales en la interfaz, y una fase final orientada a revisión
con usuarios, estabilidad e internacionalización. La distribución temporal y el
grado de cumplimiento de estas iteraciones se detallan en la
Tabla~\ref{tabla:sprints_cronologia}.

\tablaSmall{Cronología de sprints y grado de ejecución del proyecto}
{p{0.08\textwidth} p{0.14\textwidth} p{0.56\textwidth} p{0.08\textwidth} p{0.10\textwidth}}
{sprints_cronologia}
{
    \textbf{Sprint} & \textbf{Fecha fin} & \textbf{Objetivo principal} & \textbf{Issues} & \textbf{Estado} \\
}
{
    S0 & 29/10/2025  & Puesta en marcha del proyecto: configuración inicial, investigación de herramientas y arranque de documentación.  & 5/5 & 100\% \\
    S1 & 20/11/2025  & Validación y comparación de modelos de temas (LDA, Top2Vec, BERTopic, FASTopic).                                  & 6/6 & 100\% \\
    S2 & 04/12/2025  & Hiperparametrización y primeras visualizaciones de resultados experimentales.                                     & 6/6 & 100\% \\
    S3 & 17/12/2025  & Construcción del conjunto de datos base con \textit{issues} y \textit{readmes} de repositorios seleccionados.               & 4/4 & 100\% \\
    S4 & 30/12/2025  & Extensión de la ingesta a memorias en \LaTeX{} y puesta en marcha del esqueleto de \textit{processing}.           & 4/4 & 100\% \\
    S5 & 15/01/2026  & Implementación del módulo de procesamiento y consolidación del flujo analítico.                                   & 7/7 & 100\% \\
    S6 & 29/01/2026  & Conversión a aplicación desplegable, mejoras de experiencia de usuario y extensión funcional.                     & 3/3 & 100\% \\
    S7 & 19/02/2026  & Mejora de usabilidad en la página de análisis académico de \acrshort{tfg}.                                        & 5/5 & 100\% \\
    S8 & 05/03/2026  & Incorporación del \textit{datasource} de resúmenes académicos (\textit{abstracts}).                               & 6/6 & 100\% \\
    S9 & 26/03/2026  & Mejora visual de la aplicación web y refinamiento de la presentación de resultados.                               & 7/7 & 100\% \\
    S10 & 20/04/2026 & Revisiones con usuarios y ajustes derivados de pruebas funcionales.                                               & 4/4 & 100\% \\
    S11 & 04/05/2026 & Refinamiento UX, datos demo para validación, estabilización de plataforma e internacionalización.                 & 6/6 & 100\% \\
    S12 & 21/05/2026 & Usabilidad, Documentación y Despliegue.                                                                           & 6/6 & 100\% \\
    S13 & 03/06/2026 & Revisión documentación y despliegue.                                                                           & 4/4 & 100\% \\
}

El histórico de sprints refleja una planificación adaptativa coherente con la
naturaleza experimental del proyecto. Aunque algunas tareas de evaluación y
visualización se abordaron en fases tempranas, esta decisión permitió detectar
limitaciones de forma anticipada y reorientar con mayor precisión el diseño de
ingesta, procesamiento y experiencia de usuario en iteraciones posteriores.

\section{Análisis de riesgos}
Un pilar fundamental en la planificación de cualquier proyecto de ingeniería es
la identificación temprana de riesgos y la definición de planes de mitigación.
Durante las fases iniciales se identificaron los siguientes riesgos
principales:

\begin{itemize}
    \item \textbf{Riesgo 1: Limitaciones de la \acrshort{api} de GitHub.} Al depender de una fuente externa para la ingesta de datos, existía el riesgo de alcanzar el límite de peticiones (Rate Limit). \textit{Mitigación:} Se implementó el uso de tokens autenticados, peticiones paginadas y tiempos de retardo controlados (\textit{sleeps}) para garantizar una extracción continua sin bloqueos.
    \item \textbf{Riesgo 2: Altos tiempos de entrenamiento \acrshort{nlp}.} El procesamiento masivo de texto y el entrenamiento concurrente de cuatro modelos de temas (\acrshort{lda}, BERTopic, etc.) suponía un cuello de botella computacional. \textit{Mitigación:} Se optó por utilizar archivos Apache Parquet para agilizar las operaciones de \ACRshort{io} en disco y se aprovechó la aceleración por GPU local para la instanciación de los \textit{embeddings} de los Transformers.
    \item \textbf{Riesgo 3: Incompatibilidad de dependencias de Python.} Las librerías de \acrshort{ai} suelen presentar conflictos estrictos de versiones. \textit{Mitigación:} Se abordó mediante la contención temprana, aislando el entorno de desarrollo en contenedores Docker y especificando versiones exactas en un archivo \texttt{requirements.txt}.
\end{itemize}

\section{Estudio de viabilidad}
Todo proyecto de ingeniería de software requiere de un análisis previo de
viabilidad para garantizar que su ejecución sea realista en términos de
recursos, coste y adecuación al marco normativo vigente.

\subsection{Viabilidad económica}
Al tratarse de un \acrshort{tfg} de naturaleza académica, no existe un
presupuesto de inversión comercial real. No obstante, es imperativo cuantificar
el valor del desarrollo para estimar cuál habría sido su coste total si este
proyecto se hubiese encargado dentro de un entorno empresarial. El análisis
contempla tres partidas principales:

\begin{itemize}
    \item \textbf{Costes de personal:} La carga lectiva de un \acrshort{tfg} estándar de 12 créditos ECTS equivale habitualmente a unas 300 horas de dedicación. Para estimar el coste salarial real en el mercado español actual, se toma como referencia el salario bruto anual medio de un Desarrollador de Software / Ingeniero de Datos de nivel medio (\textit{Mid-level} o \textit{Semi-senior}) bajo las tablas reales de mercado y el \textit{XIX Convenio Colectivo Estatal de Empresas de Consultoría y Estudios de Mercado}~\cite{BoeConvenioTIC} (la media de mercado real para perfiles tecnológicos con experiencia intermedia y conocimientos de analítica de datos e inteligencia artificial ronda los 38.000 € brutos anuales, según datos de la \textit{Guía del Mercado Laboral 2026} de Hays~\cite{HaysGuia2026}). A esta cifra se le añade aproximadamente un 32\% en concepto de cotizaciones a la Seguridad Social a cargo del empleador (contingencias comunes, desempleo, FOGASA, formación profesional y MEI, de acuerdo con la regulación de tipos de cotización vigente~\cite{BoeCotizaciones2026}), resultando en un coste anual para la empresa de 50.160 €. Considerando la jornada anual estándar de 1.800 horas estipulada en dicho convenio, el coste por hora del desarrollador asciende a aproximadamente 27,87 €/hora. Por lo tanto, el coste total de personal por las 300 horas de dedicación estimadas se cifra en \textbf{8.360 €}.
    \item \textbf{Costes de revisión y auditoría (Tutoría):} El seguimiento, asesoramiento técnico, validación y control de calidad metodológica han sido desempeñados por el tutor del proyecto (ejerciendo funciones homólogas a las de un Consultor Senior o Director Técnico). Se estima una dedicación acumulada de 25 horas para reuniones de seguimiento, revisión de entregables y evaluación final del proyecto. Con un coste de mercado medio para este perfil de consultoría senior fijado en 45 €/hora, el coste total de esta partida asciende a \textbf{1.125 €}.
    \item \textbf{Costes de hardware (Amortización):} El desarrollo de las herramientas y el entrenamiento de los algoritmos de modelado requieren alta capacidad de cómputo. Para ello, se ha utilizado un equipo propio (AMD Ryzen 5 7600X, 32 GB RAM, NVIDIA RTX 4070 SUPER 12GB). Estimando un valor de compra de 1.800 € y una vida útil de 5 años, la amortización correspondiente a los 8 meses del proyecto supone un impacto de \textbf{240 €}.
    \item \textbf{Costes de software e infraestructura:} Todos los lenguajes, librerías, y entornos de desarrollo utilizados (Python, Visual Studio Code, Docker, \href{https://github.com/zcc1001/tfg-topics}{GitHub}) disponen de licencias de código abierto o planes gratuitos. En consecuencia, el coste en software es de \textbf{0 €}.
    \item \textbf{Gastos generales prorrateados:} Contemplan costes indirectos del proyecto a lo largo de sus 8 meses de duración, desglosados de forma mensual: alquiler proporcional del espacio de oficina o puesto de trabajo (80 €/mes), consumo eléctrico (40 €/mes), conexión a internet de banda ancha (30 €/mes) y material de oficina (10 €/mes). Esto supone un gasto general mensual de 160 €/mes, resultando en un total prorrateado de \textbf{1.280 €}.
\end{itemize}

A continuación, en la Tabla~\ref{tabla:presupuesto}, se muestra de forma
resumida el desglose presupuestario del proyecto.

\tablaSmall{Tabla resumen del presupuesto estimado del proyecto}
{l r}
{presupuesto}
{
    \textbf{Concepto}                               & \textbf{Coste estimado} \\
}
{
    Costes de personal (desarrollo, 300h)           & 8.360 €                 \\
    Revisión y auditoría (tutoría, 25h)             & 1.125 €                 \\
    Amortización de hardware                        & 240 €                   \\
    Licencias de software                           & 0 €                     \\
    Gastos generales                                & 1.280 €                 \\
    \textbf{Presupuesto Total}                      & \textbf{11.005 €}       \\
}

Agregando todas las partidas, el presupuesto total de desarrollo de esta
herramienta rondaría los \textbf{11.005 €}. Teniendo en cuenta la nula barrera
de entrada a nivel de licencias y herramientas, se concluye que el proyecto es
totalmente \textbf{económicamente viable}.

\subsection{Viabilidad legal}
Desde la perspectiva jurídica, el desarrollo, almacenamiento y futura
distribución de esta aplicación en código abierto plantea dos retos principales
que han sido solventados:

\begin{itemize}
    \item \textbf{Protección de datos personales:} Dado que el sistema requiere analizar metadatos e incidencias de repositorios que actúan como \acrshort{tfg}, existía el riesgo de vulnerar la privacidad de autores y tutores. Para solucionarlo, el módulo de extracción se ha configurado para \textbf{anonimizar intencionalmente} cualquier nombre o dato de contacto recopilado. Además, la información es extraída de forma no invasiva \textbf{exclusivamente desde repositorios públicos} de \href{https://github.com}{GitHub}, acatando escrupulosamente los Términos de Servicio de su \acrshort{api} para usos académicos.
    \item \textbf{Licenciamiento del software:} El código desarrollado hace uso de librerías de terceros (tales como Pandas o Streamlit) amparadas bajo licencias permisivas (BSD o Apache 2.0). Para asegurar la mayor interoperabilidad y evitar bloqueos en usos posteriores por otros investigadores, el código fuente completo de este proyecto será liberado en \href{https://github.com/zcc1001/tfg-topics}{GitHub} bajo la \textbf{Licencia Apache 2.0}. Esta licencia de código abierto y altamente permisiva permite a cualquier usuario usar, modificar y distribuir el software para cualquier propósito, tanto comercial como académico, garantizando al mismo tiempo la protección frente a reclamaciones de patentes y requiriendo la conservación del aviso de derechos de autor original y del descargo de responsabilidad, lo que consolida por completo su \textbf{viabilidad legal}.
\end{itemize}

